% Source: ACM_Templates.pdf, PDF pp. 81-82.
\subsection{树上点差分与边差分}
\label{tree:difference}
\AlgorithmInfo
    {离线处理多次路径加，最后统一求每个顶点或每条边的总贡献。}
    {倍增预处理 $O(n\log n)$，每次路径修改 $O(\log n)$，一次汇总 $O(n)$；总空间 $O(n\log n)$。}
    {先复制 TreeLCA；点从 1 开始，根的父亲为 0；差分中间值与汇总值均须在有符号 64 位内。}
\subsubsection*{原理}
令 $p=LCA(u,v)$，希望路径上增加 $w$。点差分执行
\[
    d_u\mathrel{+}=w,\quad d_v\mathrel{+}=w,\quad
    d_p\mathrel{-}=w,\quad d_{fa_p}\mathrel{-}=w.
\]
递归回溯时把儿子差分加到父亲：两端贡献向上汇总，在 LCA 处保留一次，再在其父亲处抵消。根的父亲 0 号位仅作占位，不当作答案。

边差分执行
\[
    d_u\mathrel{+}=w,\quad d_v\mathrel{+}=w,\quad d_p\mathrel{-}=2w.
\]
汇总后的 $d_x$ 表示边 $(fa_x,x)$ 的贡献，因为两端在该边下方子树内恰有一个时才会留下净值。两种差分分别维护，不能混用。$u=v$ 时，点差分修改一个点，边差分不修改任何边。
\subsubsection*{使用教程}
先 \InlineCode{init(n)}，用 \InlineCode{add(u,v)} 加无向边（每边一次），再 \InlineCode{build(root)} 预处理，省略根取 1。

\InlineCode{add_point(u,v,w)} 给路径含两端的每个点加 $w$，\InlineCode{add_edge(u,v,w)} 给路径每条边加 $w$。\InlineCode{collect()} 返回两份从 1 开始的数组：点贡献，以及按较深端点存放的边贡献，根对应的边结果不用。汇总使用副本，重复调用不会重复累加；\InlineCode{clear()} 清除修改并保留树和倍增。
\begin{lstlisting}[style=mtycpp,numbers=none]
TreeDifference g;
g.init(4);
g.add(1,2),g.add(2,3),g.add(2,4);
g.build(1);
g.add_point(3,4,2);
g.add_edge(1,3,5);
auto [point,edge]=g.collect();
long long vertex=point[2],incoming_edge=edge[3];
\end{lstlisting}
这是离线汇总接口；每次修改后若都调用 \InlineCode{collect()}，每次都会花费 $O(n)$，不能据此得到在线高效路径查询。
\needspace{7\baselineskip}
\subsubsection*{模板代码}
\CodeFile{code/trees/difference.cpp}
